% =====================================================================
%  CAPÍTULO 2 — Ejemplo: metodología (ecuaciones y referencias cruzadas)
% =====================================================================
\chapter{Título del capítulo 2}\label{cap:capitulo2}

\section{Introducción del capítulo}

\lipsum[9]

\section{Modelo teórico}

% Ecuación numerada: \begin{equation} ... \label{eq:...} \end{equation}
Siguiendo a \textcite{mankiw1992}, la función de producción es
\begin{equation}
  Y_t = K_t^{\alpha} H_t^{\beta} (A_t L_t)^{1-\alpha-\beta},
  \qquad \alpha + \beta < 1. \label{eq:produccion}
\end{equation}

% Varias ecuaciones alineadas en el signo "=": entorno align
El problema del hogar y su Lagrangiano son:
\begin{align}
  \max_{c,\,\ell} \; & u(c, \ell) \quad \text{s.a.} \quad c = w(1-\ell) + y, \label{eq:problema}\\
  \mathcal{L} &= u(c,\ell) + \lambda\,[w(1-\ell) + y - c]. \label{eq:lagrangiano}
\end{align}

\section{Estrategia empírica}

El modelo econométrico es
\begin{equation}
  y_i = \beta_0 + \beta_1 x_i + \mathbf{z}_i^{\top}\boldsymbol{\gamma}
        + \varepsilon_i, \qquad \E[\varepsilon_i \mid x_i, \mathbf{z}_i] = 0,
  \label{eq:modelo}
\end{equation}
que se estima por mínimos cuadrados ordinarios
\parencite{wooldridge2010, angrist2009}. La \cref{eq:modelo} se deriva de la
\cref{eq:produccion}; los resultados se presentan en el \cref{cap:capitulo3}.

\section{Datos}

\lipsum[10]
